% define \title (only used by writelatex.com)
%\title{Labbrapportsmall fysikalisk kemi}
\documentclass[a4paper,12pt]{article}

% ----------------------------------------------------------------------
% Font set up
% ----------------------------------------------------------------------
\usepackage[T1]{fontenc}
\usepackage[utf8]{inputenc}
\usepackage{lmodern}
\usepackage{mathpazo}

% ----------------------------------------------------------------------
% Support packages
% ----------------------------------------------------------------------
\usepackage[swedish]{babel}
\usepackage[margin = 1in]{geometry}
\usepackage{fancyhdr}
\usepackage{lastpage}
\usepackage{amssymb}
\usepackage{amsmath}
\usepackage{graphicx}
\usepackage{caption,subcaption}
\usepackage{booktabs}
\usepackage{multirow}
\usepackage{csquotes}
\usepackage{xcolor}
   \definecolor{maroon}{cmyk}{0.00,0.87,0.68,0.32}
   \definecolor{sepia} {cmyk}{0.00,0.83,1.00,0.70}

% ----------------------------------------------------------------------
% Chemistry and units packages
% ----------------------------------------------------------------------
\usepackage[version=3]{mhchem}
\usepackage{siunitx}
   \sisetup{load-configurations = abbreviations}
   \DeclareSIUnit{\molar}{\mole\per\cubic\deci\metre}
   \DeclareSIUnit{\ppm}{ppm}
   
   
% ----------------------------------------------------------------------
% FYLL I FÖRFATTARENS UPPGIFTER
% ----------------------------------------------------------------------
\newcommand{\forfattare}{Förnamn Efternamn}
\newcommand{\program}{W2A Miljö- och vattenteknik}
\newcommand{\email}{fornamn.efternamn@student.uu.se}
% ----------------------------------------------------------------------
% FYLL I RAPPORTENS UPPGIFTER
% ----------------------------------------------------------------------
\newcommand{\titel}{Spektrofotometrisk bestämning av pKa för BTB}
\newcommand{\rapportnummer}{Laboration 2}
\newcommand{\datum}{2013-11-18}
\newcommand{\medlaboranter}{Förnamn Efternamn, Förnamn Efternamn}
% ----------------------------------------------------------------------
% KURSENS UPPGIFTER
% ----------------------------------------------------------------------
\newcommand{\labass}{Förnamn Efternamn}
\newcommand{\kurs}{Fysikalisk kemi, 1KB301}


% ----------------------------------------------------------------------
% setup pdf meta-information
% ----------------------------------------------------------------------
\usepackage[pdfkeywords={},%
            pdfauthor={\forfattare, \program},%
            pdftitle={\titel},%
	        pdfsubject={\kurs},%
	        colorlinks=true,%
            linkcolor=sepia,%
	        citecolor=sepia,%
	        filecolor=maroon,%
	        urlcolor=sepia]{hyperref}
\usepackage{cleveref}
\input{swedishcref.tex}




% ----------------------------------------------------------------------
% newcommands...
% ----------------------------------------------------------------------
\newenvironment{fraga}[1]{\subsection{#1}}{}


\begin{document}

% ----------------------------------------------------------------------
% Försättsblad
% ----------------------------------------------------------------------
\begin{titlepage}
\thispagestyle{fancy}
\fancyhf{}
\setlength{\headheight}{28pt}
\renewcommand{\headrulewidth}{1ex}
\renewcommand{\footrulewidth}{1ex}
\fancyfoot[R]{\texttt{\email}}
\fancyhead[R]{Uppsala universitet}
\fancyhead[L]{\forfattare\\\program}
\begin{flushleft}%
\texttt{Laboration utförd: \datum}\\
\texttt{Medlaboranter: \medlaboranter}\\
\texttt{Lab.\thinspace ass.: \labass}\\
\texttt{\kurs}
\end{flushleft}%
\vspace{\stretch{8}}
\begin{center}
{\bf\huge{\parbox{0.8\linewidth}{\centering\titel}}}
\end{center}
\vspace{\stretch{1}}
\begin{center}
\texttt{\rapportnummer}\\
\texttt{Ångströmslaboratoriet}\\
\texttt{\today}%\\
\end{center}
\vspace{\stretch{15}}
\end{titlepage}


% ----------------------------------------------------------------------
% Innehållsförteckning
% ----------------------------------------------------------------------
\thispagestyle{empty}
\tableofcontents
\newpage

% ----------------------------------------------------------------------
% återställ sidfot och sidhuvud
% ----------------------------------------------------------------------
\setcounter{page}{1}
\pagestyle{fancy}
\setlength{\headheight}{15pt}
\fancyhf{}
\fancyhead[R]{\forfattare}
\fancyfoot[R]{\thepage\ (\pageref{LastPage})}




\section{Introduktion}
Plancks lag kan användas för att beskriva strålningen från en svartkropp, $I$, som funktion av våglängd, $\lambda$, och yttemperatur, $T$:
\begin{equation}
\label{eq:planck}
I\left(\lambda, T\right) = \frac{2{\pi}hc^2}{\lambda^5}%
   \frac{1}{\exp\left(\displaystyle\frac{hc}{{\lambda}kT}\right) - 1}
\end{equation}
Storheten $I$ har enheten $[I]=\si{\watt\per\cubic\metre}$, vilket vi kan bekräfta med en enhetsanalys:
\begin{equation}
\label{eq:planck-unit}
\left[I(\lambda, T)\right] = \left[\frac{hc^2}{\lambda^5}\right] =%
   \frac{\si{\joule}\times\si{\second}\times\left(\si{\metre}\times\si{\per\second}\right)^2}{\si{\metre}^5} =%
   \si{\joule\per\second\per\cubic\metre} =% 
   \si{\watt\per\cubic\metre}
\end{equation}
Som vi såg i \cref{eq:planck-unit}, har $I$ enheten \si{\watt\per\cubic\metre}.

\ce{TiCl4} bildar i fuktig luft hydroxiden \ce{Ti(OH)4} under avgivande av saltsyra.
\begin{equation}
\ce{TiCl4(l) + 4H2O(g) -> Ti(OH)4(s) + 4HCl(g)}
\end{equation}
Hydroxiden avger vatten och bildar den termodynamiskt stabila fasen \ce{TiO2}.
\begin{equation}
\ce{Ti(OH)4 ->[][-H2O] TiO2.H2O ->[][-H2O] TiO2 \text{partiklar}}
\end{equation}

Följande är ett exempel på en kemisk jämviktsreaktion:
\begin{equation}
\ce{H2 + I2  <=> 2HI}
\end{equation}
vars jämviktskonstant kan skrivas:
\begin{equation}
K = \frac{\ce{\{HI\}}^2}{\ce{\{H2\}}\cdot\ce{\{I2\}}}
\end{equation}

För en enprotonig syra, det vill säga för en syra \ce{HA} som dissocierar i endast ett steg, kan vi skriva följande generella jämviktsreaktion, med jämviktskonstanten $K_\mathrm{a}$:
\begin{equation}
\ce{HA <=> H+ + A-}
\end{equation}
Av den allmänna definitionen av jämviktskonstant följer att:
\begin{equation}
K_\mathrm{a} = \frac{{\{\ce{H+}\}\cdot\{\ce{A-}\}}}{{\{\ce{HA}\}}}
\end{equation}

Buffertformeln eller Henderson-Hasselbalch-ekvationen är en enkel formel för beräkningen av pH i en buffertlösning. Den kan skrivas som
\begin{equation}
\mathrm{pH} = \mathrm{p}K_\mathrm{a} + \log_{10}\frac{[\ce{A-}]}{[\ce{HA}]}
\end{equation}
där [\ce{A-}] är koncentrationen av buffertens basform och [\ce{HA}] är koncentrationen av buffertens syraform. I formeln ingår också buffertens $\mathrm{p}K_\mathrm{a}$.


\section{Utförande}
Vanlig brödtext, \emph{text med emfas}, \textbf{fetstil}.

Börja nytt stycke genom att lämna en tom rad.

Exempel på en onumrerad lista:
\begin{itemize}
\item $\alpha$
\item $\beta$
\item $\gamma$
\end{itemize}


\subsection{Beredning av buffertlösningar}
\SI{2.00e-4}{\molar} \ce{BTB}, \SI{0.100}{\molar} \ce{HCl}, \SI{0.100}{\molar} \ce{NaOH}, samt \ce{KH2PO4(aq)}.

\begin{figure}[tbh]
\centering
\begin{subfigure}[b]{0.48\textwidth}
\centering
\includegraphics[width=\linewidth]{Bromothymol_blue_colors.jpg}
\caption{Sur, neutral och basisk BTB-lösning.}
\label{fig:BTB-losningar}
\end{subfigure}
\,%
\begin{subfigure}[b]{0.48\textwidth}
\centering
\includegraphics[width=\linewidth]{Bromothymol-blue-2D-skeletal.png}
\caption{Strukturformel \ce{HBTB} (protonerad).}
\label{fig:BTB-struktur}
\end{subfigure}%
\caption{Bromtymolblått (BTB) är ett färgämne som ofta används för att i en vattenlösning indikera surhetsgrad, pH, där man utnyttjar att bromtymolblått byter färg vid neutralt pH. Bromtymolblåtts gul-orangea komplex dominerar vid pH under 7 (surt) och dess blåa vid pH över 7 (basiskt). Vid neutralt är koncentrationerna av de båda komplexen lika stora vilket resulterar i en grön färg.}
\label{fig:BTB}
\end{figure}



\subsection{Spektrofotometrisk bestämning}
Text


\section{Resultat}
Text

\begin{table}[tbh]
\centering
\caption{Exempel på en tabell.}
\label{tab:exempel}
\begin{tabular}{S[table-format=1.2]S[table-format=1.3]S[table-format=1.2e+2]S[table-format=1.2e+2]}
\toprule
{pH} & {Absorbans} & {\ce{[BTB^-]}/\si{\molar}} &  {\ce{[HBTB]}/\si{\molar}}\\\midrule
6.03 & 0.098 & 2.41e-6 & 3.64e-5\\
6.54 & 0.423 & 1.04e-5 & 5.01e-5\\
7.03 & 1.048 & 2.58e-5 & 3.47e-5\\
\bottomrule
\end{tabular}
\end{table}

\begin{figure}[tbh]
\centering
\includegraphics[width=\textwidth]{Campus_A_204H.jpg}
\caption{Exempel på en bild.}
\label{fig:exempel}
\end{figure}


\section{Diskussion och slutsatser}
Text


\section{Svar på obligatoriska frågor}

\begin{fraga}{Fråga 1}
Svar på fråga 1
\end{fraga}

\begin{fraga}{Fråga 2}
Svar på fråga 2
\end{fraga}






\newpage
\section{Bibliografi}
\begin{enumerate}
\item Shriver, D.~F.~\&~Atkins, P.~W.~(1999), \textit{Inorganic Chemistry, third ed}. Oxford: Oxford University Press.
\item Wikipedia-bidragsgivare, \textit{Bromtymolblått}, Wikipedia, \url{http://sv.wikipedia.org/w/index.php?title=Bromtymolbl%C3%A5tt&oldid=20117267} (hämtad 16 nov 2013).
\item Wikipedia-bidragsgivare, \textit{Buffertformeln}, Wikipedia, \url{http://sv.wikipedia.org/w/index.php?title=Buffertformeln&oldid=20378186} (hämtad 16 nov 2013).
\item Wikipedia-bidragsgivare, \textit{Syrakonstant}, Wikipedia, \url{http://sv.wikipedia.org/w/index.php?title=Syrakonstant&oldid=20170909} (hämtad 16 nov 2013).
\item Wikipedia-bidragsgivare, \textit{Kemisk jämvikt}, Wikipedia, \url{http://sv.wikipedia.org/w/index.php?title=Kemisk_j%C3%A4mvikt&oldid=20161143} (hämtad 16 nov 2013). 
\end{enumerate}



\end{document}